\documentclass[../../thesis_main.tex]{subfiles}
\externaldocument{../../thesis_main}
\graphicspath{{\subfix{../../figures/}}}

\begin{document}
    % Elenchi delle note: per stato e per livello
    \ifSubfilesClassLoaded{\notesummary}{}
    \ifSubfilesClassLoaded{\listofthesisnotes}{}
    \ifSubfilesClassLoaded{\listofgeneralnotes}{}
    \ifSubfilesClassLoaded{\listofopennotes}{}
    \ifSubfilesClassLoaded{\listofsolvednotes}{}
    \ifSubfilesClassLoaded{\listoftodos}{}
    \ifSubfilesClassLoaded{\listofrevnotes}{}

    \ifSubfilesClassLoaded{\tableofcontents}{}

    \chapter{Theoretical Background}
This chapter provides an overview of the theoretical background underlying the main topics addressed in this thesis.\CCgen{capitolo 1}{la notazione non e' unificata fra le sezioni}{Tre simboli cambiano significato passando da una sezione all'altra, e sono contati sul sorgente. $\nu$ e' la frequenza ottica dei fasci nella sezione sul conveyor belt (32 occorrenze) e la frequenza di trappola nella sezione EIT (22). $\Omega$ e' la frequenza di trappola nella prima (3 occorrenze) e la frequenza di Rabi nella seconda (15). La larghezza naturale e' $\Gamma$ nella sezione sulle trappole (20 occorrenze) e $\gamma$ nella sezione EIT (20): stessa grandezza fisica, due simboli, e in nessuno dei due punti si dice che sono la stessa cosa. Nel capitolo 3 la tabella dei parametri usa $\nu_z$ e $\gamma$ insieme, quindi il lettore deve tenere a mente entrambe le convenzioni. Due modi di sistemarla: una tabella dei simboli in apertura di capitolo, oppure rinominare (per esempio $\nu_{\mathrm{trap}}$ per la trappola e un solo simbolo per la larghezza in tutta la tesi).}\CCnote{come leggere questo PDF}{\textbf{Come leggere questo PDF.} Due colonne di note. A SINISTRA quelle di MD, invariate, con la loro lista in testa al capitolo. A DESTRA quelle degli altri revisori, marcate nel testo da una sottolineatura tratteggiata e da un numero, per esempio [CC 12], che compare identico sul riquadro: e' quello il legame fra nota e frase, anche quando i riquadri scivolano in basso. \textbf{Aggiornato a \texttt{c977f7c}.} Rispetto al giro precedente hai risolto 16 rilievi, e le note corrispondenti sono state tolte: fra questi il termine $\delta_{2\mathrm{ph}}$ mancante in $H_{\mathrm{rot}}$, la costante $E_{g_1}$, la diagonale della matrice 2x2 con le etichette $\pm$ che ne dipendevano, le decine/centinaia di THz, il rimando sbagliato alla polarizzabilita', le due figure del conveyor mai richiamate, i credits della figura dei modi, e le citazioni Benabid 2002 e Joannopoulos. Restano aperte le altre, piu' tre rilievi nuovi introdotti da questa revisione.}

\CCgen{capitolo 1}{distinguere il background del lavoro fatto da quello del lavoro futuro}{Il capitolo tratta allo stesso modo cose molto diverse quanto a ruolo nella tesi. La melassa e l'ODT sono lo sfondo di misure effettivamente fatte; il conveyor belt e' derivato per intero (25 equazioni) ma non e' mai stato realizzato, e il raffreddamento EIT e' simulato, non eseguito. Il lettore non ha modo di capirlo dal capitolo 1: tutto e' presentato con lo stesso tono. Una frase in apertura che dica quali sezioni preparano risultati sperimentali e quali preparano il lavoro futuro orienta la lettura e, soprattutto, mette al riparo dall'impressione che il conveyor belt sia stato messo in funzione.} The first part introduces hollow-core fibers and describes their light-guiding mechanisms. This is followed by an overview of atomic trapping techniques, starting from the magneto-optical trap(MOT) and the \MDssolved{optical dipole trap}{aggiungerei "(ODT)" per coerenza stilistica}{Fatto: il testo ora scrive ``the optical dipole trap(ODT)'', quindi la sigla e' introdotta qui come lo e' quella del MOT poco prima. Resta da sistemare una minuzia tipografica che l'inserimento si e' portato dietro, e vale per entrambe: manca lo spazio prima della parentesi, si legge \texttt{trap(MOT)} e \texttt{trap(ODT)} invece di \texttt{trap (MOT)} e \texttt{trap (ODT)}.}(ODT), outlining the fundamental principles that enable the cooling and confinement of neutral atoms.

\MDs{Building}{"Building" non suona molto bene. Io invertirei la frase per renderla più liscia} on these concepts, more advanced manipulation techniques are then discussed. Specifically, the concept of the optical conveyor belt is introduced, demonstrating how moving optical lattices \MDs{allow for the controlled transport of atoms}{non mi sembra inglese corretto. Io direi "enables the control of the transport of atoms"}. Finally, the theoretical principles of electromagnetically induced transparency (EIT) cooling are presented, highlighting its potential for achieving sub-Doppler temperatures inside the fiber core.
    \import{./}{cap1_HCPCF.tex}
    \clearpage
    \import{./}{cap1_trapping.tex}
    \clearpage
     \import{./}{cap1_Conveyor_belt.tex}
    \clearpage
     \import{./}{cap1_EIT.tex}
     
    \ifSubfilesClassLoaded{\printbibliography}{}
\end{document}